\chapter{Bourses, courtiers et plateformes}\label{ch:m1:exchanges-brokers-venues}

L'entreprise qui possède le New York Stock Exchange a tiré en 2025 davantage
de la vente de données et de connexions réseau dans ses \omterm{def:m1:exchanges-brokers-venues:exchange}{bourses} que de
l'ensemble des actions et des options sur actions qui y étaient négociées. Le fameux parquet
ne représente qu'une petite partie d'une activité qui cote des sociétés, apparie des ordres, compense des transactions, loue des baies
dans un centre de données et concède sous licence l'enregistrement de ce qui s'est passé. Une société de trading est
cliente de ces cinq services --- et ne peut envoyer un seul ordre avant de décider
par l'intermédiaire de qui, et sous quel nom, il atteint le moteur d'appariement. Ce
chapitre décrit ce que vend une \omterm{def:m1:exchanges-brokers-venues:exchange}{bourse}, la chaîne d'intermédiaires
entre un trader et l'appariement, ainsi que la multitude de plateformes qui ne sont pas
des \omterm{def:m1:exchanges-brokers-venues:exchange}{bourses}.

\section{Ce que vend une bourse}

\begin{definition}[Ordre]\label{def:m1:exchanges-brokers-venues:order}
Un \emph{ordre}\index{ordre} est une instruction ferme d'achat ou de vente
d'une quantité déterminée d'un instrument, comportant au minimum un sens, une quantité,
une condition de prix (une limite, ou «~au marché~») et une durée de validité. Il reste de la
responsabilité de son émetteur jusqu'à ce qu'il soit exécuté, annulé ou expiré.
\end{definition}

\begin{definition}[Bourse]\label{def:m1:exchanges-brokers-venues:exchange}
Une \emph{bourse}\index{bourse} est une organisation réglementée qui admet
des instruments à la négociation, admet des membres et exploite un système dans lequel
les ordres de ses membres interagissent selon des règles publiques et non discrétionnaires pour former
des transactions et des prix. Elle supervise son propre marché et est elle-même supervisée par
une autorité publique.
\end{definition}

La définition recouvre cinq produits, vendus à des clients différents.

\begin{method}[Lire le chiffre d'affaires d'un groupe de bourses]\label{met:m1:exchanges-brokers-venues:products}
Classer chaque ligne de chiffre d'affaires en~:
\begin{enumerate}
\item \textbf{Cotation} --- payée par les émetteurs, annuellement, pour l'admission
  à la négociation.
\item \textbf{Transaction} --- payée par les membres par action, contrat ou unité de valeur
  appariée, nette des éventuelles ristournes qui leur sont reversées.
\item \textbf{Compensation} --- lorsque le groupe possède la chambre de
  compensation (\cref{ch:m1:clearing-and-settlement}), payée par transaction compensée et
  perçue sur le collatéral détenu.
\item \textbf{Données} --- licences pour les cours temps réel et historiques, payées par
  tous ceux qui les consultent ou les utilisent pour calculer.
\item \textbf{Accès} --- ports, interconnexions, baies et alimentation électrique dans le
  centre de données de la \omterm{def:m1:exchanges-brokers-venues:exchange}{bourse}, payés par ceux qui doivent se trouver à proximité
  (livre~14).
\end{enumerate}
Les lignes 1, 4 et 5 se répètent quels que soient les volumes~; les lignes 2 et 3 suivent le marché.
Il faut diviser les lignes 2 et 3 par le volume pour obtenir le \emph{chiffre d'affaires par contrat} ou
par action~: c'est le prix du service central de la \omterm{def:m1:exchanges-brokers-venues:exchange}{bourse}.
\end{method}

\begin{dated}{2026-09}{Deux groupes boursiers à travers leurs propres chiffres}\label{dat:m1:exchanges-brokers-venues:groups}
\textbf{Intercontinental Exchange}, propriétaire du New York Stock Exchange, a publié
un chiffre d'affaires net de 9.9~milliards de dollars pour 2025, dont 5\,411~millions de
dollars dans son segment \omterm{def:m1:exchanges-brokers-venues:exchange}{Bourses}~: énergie 2\,182~millions de dollars, produits agricoles et métaux
233~millions de dollars, contrats à terme financiers 608~millions de dollars, actions au comptant et
options sur actions 467~millions de dollars, gré à gré et autres 395~millions de dollars, services de données et
de connectivité 1\,031~millions de dollars, cotations 495~millions de dollars.
\textbf{CME Group} a publié un chiffre d'affaires de 6.5~milliards de dollars pour 2025, dont
5\,281~millions de dollars en frais de compensation et de transaction et 803~millions de dollars
en données de marché~; son tarif moyen par contrat était de 0.696~\$.
\end{dated}

\begin{omfigure}
\begin{tikzpicture}
\begin{axis}[width=11.5cm,height=5.2cm,xbar,bar width=7pt,
  xlabel={Chiffre d'affaires net 2025 (millions de \$)},
  ytick=data,yticklabels={{Énergie},{Agricoles et métaux},{Produits financiers},{Actions et options sur actions},{OTC et autres},{Données et connectivité},{Admissions à la cote}},
  yticklabel style={font=\small},tick label style={font=\small},label style={font=\small},
  y dir=reverse,xmin=0,xmax=2500,xtick={0,500,1000,1500,2000,2500},scaled x ticks=false,
  xticklabel style={/pgf/number format/1000 sep={\,}},enlarge y limits=0.1,grid=major,grid style={gray!25},
  table/col sep=comma]
\addplot[fill=omDef!60,draw=omDef] table[x=usd_m,y=k]{figdata/markets-1/04-exchanges-brokers-venues/ice_exchanges_2025.csv};
\end{axis}
\end{tikzpicture}

\omcaption{Le segment \omterm{def:m1:exchanges-brokers-venues:exchange}{Bourses} d'un groupe en 2025. Les données et les cotations ---
les deux barres du bas, 28~\% du segment --- ne dépendent pas des volumes~; la négociation
des actions représente moins d'un dixième. Source~: le communiqué de résultats du groupe, comme
dans l'\cref{dat:m1:exchanges-brokers-venues:groups}.}\label{fig:m1:exchanges-brokers-venues:ice}
\end{omfigure}

\begin{example}[Ce que rapporte un contrat à terme à sa bourse]\label{ex:m1:exchanges-brokers-venues:rpc}
À 0.696~\$ par contrat, 5\,281~millions de dollars de frais de compensation et de transaction
correspondent à environ 7.6~milliards de côtés de contrats facturés dans l'année, soit près de 30
millions par séance de négociation. Un seul contrat sur indice d'actions contrôle plusieurs
centaines de milliers de dollars d'actions~: les frais de la \omterm{def:m1:exchanges-brokers-venues:exchange}{bourse} représentent une faible fraction
d'un point de base de la valeur négociée. Les \omterm{def:m1:exchanges-brokers-venues:exchange}{bourses} sont, comme les \omterm{def:m1:what-a-trading-firm-does:market-maker}{teneurs de marché}, une
activité de chiffres minuscules multipliés par des chiffres énormes --- avec cette différence
que les clients d'une \omterm{def:m1:exchanges-brokers-venues:exchange}{bourse} ne peuvent pas facilement prendre un contrat ailleurs.
\end{example}

\begin{remark}[Pourquoi les bourses de produits à terme gagnent davantage que les bourses d'actions]\label{rem:m1:exchanges-brokers-venues:silo}
Une action est la même action sur chaque plateforme~: le droit américain et le droit européen
imposent la concurrence entre les \omterm{def:m1:exchanges-brokers-venues:exchange}{bourses} d'actions, et les frais d'appariement ont été réduits
presque à rien. Un contrat à terme est une création de la \omterm{def:m1:exchanges-brokers-venues:exchange}{bourse} qui le cote
et qui le compense dans sa propre chambre de compensation~; une position ouverte là-bas ne peut
être clôturée que là-bas. Le «~silo vertical~» formé par la cotation, l'appariement et la compensation
explique pourquoi la \cref{fig:m1:exchanges-brokers-venues:ice} a cette allure.
\end{remark}

\section{Membres, courtiers et accès}

\begin{definition}[Membre de bourse]\label{def:m1:exchanges-brokers-venues:member}
Un \emph{membre de bourse}\index{membre de bourse} (ou \emph{participant}) est une
société admise par la \omterm{def:m1:exchanges-brokers-venues:exchange}{bourse} à envoyer des ordres à son système en son propre nom.
Elle doit être agréée par un régulateur, satisfaire aux exigences de fonds propres, accepter le
règlement de marché, et disposer d'un dispositif de compensation de ses transactions.
\end{definition}

\begin{definition}[Courtier]\label{def:m1:exchanges-brokers-venues:broker}
Un \emph{courtier}\index{courtier} (broker) est un membre qui envoie des ordres pour le compte de
clients qui ne sont pas membres. La \omterm{def:m1:exchanges-brokers-venues:exchange}{bourse} ne connaît que le courtier~: envers le marché, le
courtier répond de chaque ordre, et du paiement de chaque
transaction, de chacun des clients qui se trouvent derrière lui.
\end{definition}

Cette dernière phrase explique la structure de l'industrie. Parce que le
\omterm{def:m1:exchanges-brokers-venues:broker}{courtier} répond de ses clients, il doit contrôler ce qu'ils envoient~; parce que
ces contrôles prennent du temps, les firmes les plus rapides cherchent à les éviter~; et
les régulateurs décident jusqu'où ils peuvent aller.

\begin{definition}[Accès direct au marché et accès sponsorisé]\label{def:m1:exchanges-brokers-venues:dma}
Dans le cadre de l'\emph{accès direct au marché}\index{accès direct au marché} (DMA),
un client envoie ses ordres par les systèmes du \omterm{def:m1:exchanges-brokers-venues:broker}{courtier}, qui les contrôlent et
les transmettent à la \omterm{def:m1:exchanges-brokers-venues:exchange}{bourse} sous l'agrément de membre du \omterm{def:m1:exchanges-brokers-venues:broker}{courtier}. Dans le cadre de
l'\emph{accès sponsorisé}\index{accès sponsorisé}, le client se connecte directement à la \omterm{def:m1:exchanges-brokers-venues:exchange}{bourse},
avec ses propres lignes et ses propres machines, en utilisant l'identifiant de membre
du \omterm{def:m1:exchanges-brokers-venues:broker}{courtier}~; les contrôles de risque du \omterm{def:m1:exchanges-brokers-venues:broker}{courtier} s'appliquent à ce flux sans qu'il transite
par l'infrastructure de traitement des ordres du \omterm{def:m1:exchanges-brokers-venues:broker}{courtier}.
\end{definition}

\begin{omfigure}
\begin{tikzpicture}[font=\small,
  box/.style={draw=omDef,rounded corners=2pt,minimum width=2.4cm,minimum height=0.75cm,align=center,fill=omDef!4},
  chk/.style={draw=omThm,rounded corners=2pt,minimum width=2.2cm,minimum height=0.75cm,align=center,fill=omThm!5},
  ex/.style={draw=omExo,rounded corners=2pt,minimum width=1.6cm,minimum height=5.9cm,align=center,fill=omExo!6},
  flow/.style={-{Stealth[length=2mm]},thick,omDef}]
\node[ex] (ex) at (13.2,-2.55) {Bourse};
\foreach \y/\lab in {0/Algorithme du courtier,-1.7/Accès direct au marché,-3.4/Accès sponsorisé,-5.1/Adhésion directe}{
  \node[anchor=west,font=\small\bfseries] at (-0.9,\y+0.74) {\lab};
}
\node[box] (c1) at (0.4,0) {Client};
\node[box] (b1) at (4.2,0) {Trader et algorithme\\du courtier};
\node[chk] (k1) at (8.4,0) {Contrôles de risque\\du courtier};
\draw[flow] (c1) -- (b1); \draw[flow] (b1) -- (k1); \draw[flow] (k1) -- (k1 -| ex.west);
\node[box] (c2) at (0.4,-1.7) {Algorithme\\du client};
\node[chk] (k2) at (8.4,-1.7) {Contrôles de risque\\du courtier};
\draw[flow] (c2) -- node[above,font=\footnotesize]{par les systèmes du courtier} (k2); \draw[flow] (k2) -- (k2 -| ex.west);
\node[box] (c3) at (0.4,-3.4) {Algorithme\\du client};
\node[chk] (k3) at (8.4,-3.4) {Contrôles pilotés par le courtier\\chez le client};
\draw[flow] (c3) -- node[above,font=\footnotesize]{lignes propres du client} (k3); \draw[flow] (k3) -- (k3 -| ex.west);
\node[box] (c4) at (0.4,-5.1) {Algorithme\\du membre};
\node[chk] (k4) at (8.4,-5.1) {Contrôles de risque\\propres du membre};
\draw[flow] (c4) -- (k4); \draw[flow] (k4) -- (k4 -| ex.west);
\end{tikzpicture}

\omcaption{Quatre façons d'atteindre le même moteur d'appariement, de la plus lente et la moins
coûteuse à mettre en place (en haut) à la plus rapide et la plus exigeante (en bas). Dans
chaque ligne, les contrôles de pré-négociation d'une firme régulée (en rouge) se situent
entre l'algorithme et la \omterm{def:m1:exchanges-brokers-venues:exchange}{bourse}~: ce qui change, c'est à qui ils appartiennent et où
ils s'exécutent.}\label{fig:m1:exchanges-brokers-venues:access}
\end{omfigure}

\begin{remark}[La fin de l'accès nu]\label{rem:m1:exchanges-brokers-venues:15c35}
Jusqu'en 2010, certains \omterm{def:m1:exchanges-brokers-venues:broker}{courtiers} américains prêtaient leur identifiant à des clients rapides,
sans aucun contrôle de pré-négociation (accès «~nu~»). Le 3 novembre 2010, la SEC a adopté
la Rule 15c3-5, la \emph{market access rule}~: un \omterm{def:m1:the-sell-side:sell-side}{broker-dealer} qui donne accès au marché,
ou qui le fournit, doit appliquer des contrôles de risque de pré-négociation --- seuils de
crédit et de fonds propres, rejet des ordres erronés, vérifications réglementaires --- qui
sont sous son \emph{contrôle direct et exclusif}. Les règles européennes sur le trading
algorithmique et sur l'accès électronique direct imposent l'équivalent. Les contrôles qu'une
firme met en place pour satisfaire à ces exigences sont un composant de la mini-firme de cette
série (livre~13)~; ce qu'ils doivent détecter est la leçon de l'incident de 2012
étudié dans ce livre.
\end{remark}

\begin{method}[Choisir un modèle d'accès]\label{met:m1:exchanges-brokers-venues:choose}
\begin{enumerate}
\item \textbf{Besoin de latence.} Si la stratégie se joue sur la vitesse au
  moteur d'appariement, seuls l'\omterm{def:m1:exchanges-brokers-venues:dma}{accès sponsorisé} ou l'agrément de membre conviendront.
\item \textbf{Volume.} Calculer le coût mensuel total
  $F + \varphi V$ de chaque modèle (coût fixe $F$, frais par action $\varphi$,
  volume $V$) et déterminer les volumes de bascule $V_{ab} = (F_b - F_a)/(\varphi_a
  - \varphi_b)$.
\item \textbf{Appétit réglementaire.} Devenir membre implique de devenir un \omterm{def:m1:the-sell-side:sell-side}{broker-dealer} ou une
  entreprise d'investissement régulée~: fonds propres, personnel de conformité, reporting,
  contrôles.
\item \textbf{Contrepartie.} Un \omterm{def:m1:exchanges-brokers-venues:broker}{courtier} peut couper un client en cas de crise~; un
  membre ne dépend que de sa firme de compensation (\cref{ch:m1:clearing-and-settlement}).
\end{enumerate}
Le choix est réexaminé à mesure que la firme grandit~;
\cref{ch:m1:getting-access-equities,ch:m1:getting-access-futures-options}
examinent des grilles tarifaires réelles.
\end{method}

\begin{example}[Quatre modèles, trois bascules]\label{ex:m1:exchanges-brokers-venues:crossover}
Prenons des frais illustratifs de 0.30 cent par action et sans coût fixe pour l'algorithme
d'un \omterm{def:m1:exchanges-brokers-venues:broker}{courtier}~; 0.10 cent plus 5\,000~\$ par mois pour le DMA~; 0.04 cent plus
25\,000~\$ par mois pour l'\omterm{def:m1:exchanges-brokers-venues:dma}{accès sponsorisé} (le client paie désormais ses propres lignes et
son hébergement en colocation)~; et 150\,000~\$ par mois, sans frais par action, pour l'agrément
de membre propre à une firme (conformité, fonds propres, cotisations). Les bascules sont
$5\,000/0.002 = 2.5$ millions d'actions par mois, $20\,000/0.0006 = 33$ millions,
et $125\,000/0.0004 = 313$ millions. Une firme qui échange 15 millions d'actions par
\emph{jour} a dépassé la dernière~; un fonds qui échange un million d'actions par mois ne devrait
même pas envisager le DMA (\cref{fig:m1:exchanges-brokers-venues:cost}).
\end{example}

\begin{omfigure}
\begin{tikzpicture}
\begin{axis}[width=11.5cm,height=5.4cm,xmode=log,ymode=log,
  xlabel={actions échangées par mois (millions)},ylabel={coût total (cents par action)},
  tick label style={font=\small},label style={font=\small},
  xmin=0.1,xmax=2000,ymin=0.005,ymax=20,grid=major,grid style={gray!25},
  legend columns=4,legend style={font=\footnotesize,at={(0.5,-0.32)},anchor=north,draw=none,fill=none,/tikz/every even column/.append style={column sep=6pt}},
  table/col sep=comma]
\addplot[omExo,very thick] table[x=shares_m,y=broker_algorithm]{figdata/markets-1/04-exchanges-brokers-venues/access_cost.csv};
\addplot[omMeth,very thick] table[x=shares_m,y=direct_market_access]{figdata/markets-1/04-exchanges-brokers-venues/access_cost.csv};
\addplot[omDef,very thick] table[x=shares_m,y=sponsored_access]{figdata/markets-1/04-exchanges-brokers-venues/access_cost.csv};
\addplot[omThm,very thick] table[x=shares_m,y=own_membership]{figdata/markets-1/04-exchanges-brokers-venues/access_cost.csv};
\legend{algorithme du courtier,DMA,accès sponsorisé,propre agrément}
\end{axis}
\end{tikzpicture}

\omcaption{Coût par action des quatre modèles d'accès en fonction du volume mensuel (les deux
axes sont logarithmiques). Le modèle le moins cher change trois fois sur quatre décades
de volume. Données~: les frais indicatifs
de l'\cref{ex:m1:exchanges-brokers-venues:crossover}, calculés par le
script du chapitre.}\label{fig:m1:exchanges-brokers-venues:cost}
\end{omfigure}

\section{Plateformes qui ne sont pas des bourses}

Apparier des ordres selon des règles ne requiert pas d'agrément de \omterm{def:m1:exchanges-brokers-venues:exchange}{bourse}.

\begin{definition}[Système de négociation alternatif]\label{def:m1:exchanges-brokers-venues:ats}
En droit américain, un \emph{système de négociation alternatif}\index{système de négociation alternatif}
(ATS) est un système qui met en relation acheteurs et vendeurs de titres comme le ferait une \omterm{def:m1:exchanges-brokers-venues:exchange}{bourse}, mais
est exploité par un \omterm{def:m1:the-sell-side:sell-side}{broker-dealer} en vertu de la Regulation ATS
(adoptée en décembre 1998) au lieu de s'enregistrer comme \omterm{def:m1:exchanges-brokers-venues:exchange}{bourse}. Il ne
procède pas à la cotation de titres et ne réglemente pas ses participants.
\end{definition}

\begin{definition}[Système multilatéral de négociation]\label{def:m1:exchanges-brokers-venues:mtf}
En droit européen, un \emph{système multilatéral de négociation}\index{système multilatéral de négociation}
(MTF) est un système multilatéral exploité par une entreprise d'investissement ou un opérateur de marché,
qui réunit de multiples intérêts acheteurs et vendeurs de tiers
sur des instruments financiers, selon des règles non discrétionnaires, d'une manière
qui aboutit à la conclusion d'un contrat.
\end{definition}

La plupart des \emph{dark pools} (voir le \cref{ch:m1:us-equity-market-structure}) sont des ATS~;
les grandes plateformes actions paneuropéennes qui concurrencent les \omterm{def:m1:exchanges-brokers-venues:exchange}{bourses} nationales
sont des MTF (voir le \cref{ch:m1:european-equity-market-structure}). Les deux négocient
des instruments qu'une \omterm{def:m1:exchanges-brokers-venues:exchange}{bourse} a admis à la cote, en profitant de son admission
et du travail de surveillance qu'elle assure, ce qui explique en partie les prix que les \omterm{def:m1:exchanges-brokers-venues:exchange}{bourses}
pratiquent pour leurs données.

\begin{proposition}[Part de marché au seuil de rentabilité d'une nouvelle plateforme]\label{prop:m1:exchanges-brokers-venues:breakeven}
Une plateforme dont les coûts fixes annuels s'élèvent à $C$ entre sur un marché où se négocient $M$ actions par jour
pendant $D$ jours par an. Elle conserve des frais nets $\nu$ par action appariée et perçoit
des revenus de données de marché proportionnels à sa part, $\delta$ pour 100~\%. Son
bénéfice à la part de marché $x$ est $x\,(M D \nu + \delta) - C$ et sa
part au seuil de rentabilité est
\[
x^\star \;=\; \frac{C}{M D \nu + \delta}.
\]
\end{proposition}
\begin{proof}
Le volume apparié est de $xMD$ actions par an, chacune rapportant $\nu$~; les revenus de données sont
$x\delta$~; les coûts sont fixes.
\end{proof}

\begin{example}[Cinq pour cent ou rien]\label{ex:m1:exchanges-brokers-venues:five}
Avec $M = 10$ milliards d'actions par jour, $D = 252$, $\nu = 0.03$ cent, $\delta =
50$ millions de dollars et $C = 40$ millions de dollars, $x^\star = 40/(756+50) = 5.0\%$. Avec une
part de 10~\%, la plateforme gagne 40.6~millions de dollars~; à 2~\%, elle perd 23.9~millions de dollars.
Mais la part de marché n'est pas un paramètre libre~: les traders envoient leurs ordres là où ils s'attendent à être exécutés,
c'est-à-dire là où se trouvent déjà les autres ordres. Une nouvelle plateforme doit
\emph{acheter} ses premiers pour cent --- par des remises supérieures à ses frais, par
des participations offertes aux \omterm{def:m1:exchanges-brokers-venues:broker}{courtiers} qui y acheminent leurs ordres, ou par une règle à laquelle
un groupe de traders accorde de la valeur (un ralentisseur, un carnet au point médian). La liquidité attire
la liquidité~; l'histoire des plateformes est une courte liste de survivantes.
\end{example}

\begin{omfigure}
\begin{tikzpicture}
\begin{axis}[width=10.5cm,height=4.8cm,
  xlabel={part de marché de la nouvelle plateforme (\%)},ylabel={bénéfice annuel (millions de \$)},
  tick label style={font=\small},label style={font=\small},
  xmin=0,xmax=20,ymin=-45,ymax=125,grid=major,grid style={gray!25},table/col sep=comma]
\addplot[omDef,very thick] table[x=share_pct,y=profit_m]{figdata/markets-1/04-exchanges-brokers-venues/venue_profit.csv};
\addplot[omThm,dashed] coordinates {(4.96,-45) (4.96,125)};
\node[omThm,font=\small,anchor=west] at (axis cs:5.3,100) {$x^\star = 5.0\%$};
\end{axis}
\end{tikzpicture}

\omcaption{Bénéfice d'une nouvelle plateforme en fonction de sa part de marché, avec les
paramètres de l'\cref{ex:m1:exchanges-brokers-venues:five}. Sous cinq pour cent
elle brûle la totalité de ses coûts fixes. Données~: calculées par le script du chapitre.}
\end{omfigure}

\section{La vue consolidée}

\begin{definition}[Flux de données de marché]\label{def:m1:exchanges-brokers-venues:feed}
Un \emph{flux de données de marché}\index{flux de données de marché} est le flux de messages
par lequel une plateforme publie ses ordres, ses cotations et ses transactions. Un flux
\emph{consolidé} regroupe les flux de toutes les plateformes qui négocient les mêmes
instruments en un seul meilleur cours acheteur, un seul meilleur cours vendeur et un seul fil des transactions.
\end{definition}

Dès qu'une action se négocie en vingt endroits, trois problèmes apparaissent, et chacun
fait l'objet de son propre chapitre. \emph{Quel est le prix~?} Quelqu'un doit fusionner les
flux, et la vue fusionnée est toujours un peu plus ancienne que ses sources
(voir les \cref{ch:m1:us-equity-market-structure,ch:m1:reading-market-data}).
\emph{Où l'ordre doit-il être acheminé~?} L'obligation de meilleure exécution du \omterm{def:m1:exchanges-brokers-venues:broker}{courtier} devient un
problème de routage résolu par logiciel (livre~10). \emph{Qui supervise
une manipulation qui utilise trois plateformes à la fois~?} La surveillance intermarchés est
la réponse des régulateurs, et ses détecteurs figurent dans le livre~9.

\begin{remark}[Chaque plateforme a un code]\label{rem:m1:exchanges-brokers-venues:mic}
Les systèmes identifient les plateformes par le \emph{code d'identification de marché} à quatre caractères
de la norme ISO~10383~: \texttt{XNYS} désigne le New York Stock Exchange, \texttt{XNAS}
Nasdaq. Un code \emph{d'exploitation} désigne l'organisation, et les codes de \emph{segment}
désignent ses marchés individuels~; le registre est publié chaque mois et est
gratuit. Le code, et non le nom, est ce qui figure dans les rapports de transaction, les
tables de routage et la construction de ce chapitre.
\end{remark}

\section{Tutoriel~: lire les comptes d'une bourse}\label{tut:m1:exchanges-brokers-venues:accounts}

\begin{tutorial}
\textbf{Objectif.} Transformer deux communiqués de résultats en quantités qui intéressent une
société de trading, puis calculer les points de bascule d'accès. \textbf{État final~:} les
actions cotées dans la \cref{fig:m1:exchanges-brokers-venues:ice} et les trois
points de bascule de l'\cref{ex:m1:exchanges-brokers-venues:crossover}.
\begin{tutsteps}
\item \textbf{Saisir les lignes publiées}, en indiquant leur source dans un commentaire ---
un chiffre sans source n'est qu'une rumeur.
\omcode{code/markets-1/04-exchanges-brokers-venues/python/venue_econ.py}{4}{15}{Lignes de chiffre d'affaires 2025 de deux groupes, sous forme de données.}\label{lst:m1:exchanges-brokers-venues:data}
\item \textbf{Calculer les parts.} Les lignes récurrentes (données, connectivité,
cotations) représentent $1\,526/5\,411 = 28\%$ du segment~; les actions au comptant et
les options sur actions $8.6\%$~; l'énergie $40\%$. Chez CME, les frais de compensation et de transaction
représentent $81\%$ du chiffre d'affaires et les données de marché $12\%$.
\item \textbf{Modéliser la décision d'accès.}
\omcode{code/markets-1/04-exchanges-brokers-venues/python/venue_econ.py}{23}{48}{Modèles d'accès sous forme de données, leur coût mensuel et le volume de bascule.}\label{lst:m1:exchanges-brokers-venues:access}
\item \textbf{Trouver le modèle le moins cher} à 1, 10, 100 et 1\,000 millions
d'actions par mois~: algorithme du \omterm{def:m1:exchanges-brokers-venues:broker}{courtier}, DMA, \omterm{def:m1:exchanges-brokers-venues:dma}{accès sponsorisé}, propre agrément.
\end{tutsteps}
\textbf{Pour aller plus loin.} Ajouter une colonne de latence et exclure les modèles plus lents
qu'un seuil donné~: le modèle le moins cher \emph{admissible} à 5 millions d'actions par
mois change. Ajouter ensuite les frais à paliers de volume propres à la \omterm{def:m1:exchanges-brokers-venues:exchange}{bourse}
(voir le \cref{ch:m1:getting-access-equities}), qui dépendent eux aussi du volume.
\end{tutorial}

\section{Construction~: le registre des plateformes}\label{bld:m1:exchanges-brokers-venues:registry}

\begin{build}
\bfield{Rôle} Chaque ordre, exécution, frais et message de données de marché de la
société miniature désigne une plateforme. Le registre est le seul endroit à savoir
ce qu'est chaque plateforme. \bfield{Interface} \texttt{Venue(mic, operating\_mic, name, kind, country,
currency, fee\_model)} avec \texttt{kind} parmi \texttt{exchange}, \texttt{ats},
\texttt{mtf}, \texttt{si}, \texttt{otc}~; et
\texttt{Registry.load(csv\_path)}, \texttt{get(mic)},
\texttt{by\_kind(kind)}, \texttt{segments\_of(operating\_mic)}.
\bfield{Règles} Les codes sont formés de quatre caractères alphanumériques en majuscules~; un code dupliqué, un
type inconnu ou un segment dont le code d'exploitation est absent est rejeté au moment du
chargement, et non à la première utilisation. Le registre est immuable une fois chargé.
\bfield{Tests d'acceptation} \texttt{code/firm/venues/tests/}, exécutés sur le
fichier d'exemple \texttt{data/markets-1/venues\_sample.csv}.
\bfield{Pour aller plus loin} Charger le fichier complet ISO~10383 fourni par
l'autorité d'enregistrement et indiquer combien de codes de segment actifs possède chaque code d'exploitation.
\end{build}

\begin{omsources}
\item Intercontinental Exchange, \emph{Intercontinental Exchange Reports
Strong Full Year 2025 Results}, 5 février 2026.
\item CME Group, \emph{CME Group Inc. Reports Fourth Consecutive Year of
Record Annual Revenue}, 4 février 2026.
\item US Securities and Exchange Commission, \emph{SEC Adopts New Rule
Preventing Unfiltered Market Access}, communiqué de presse 2010-210, et Release
34-63241 (Rule 15c3-5).
\item 17 CFR 242.300 (Regulation ATS, définitions)~; Directive 2014/65/EU
(MiFID II), article 4(1)(22).
\item ISO 10383 Registration Authority, \emph{Market identifier codes}.
\item L.~Harris, \emph{Trading and Exchanges}, Oxford University Press, 2003,
chapitres 3, 7 et 26.
\end{omsources}

\section{Exercices}

\begin{exercise}[$\star$]\label{exo:m1:exchanges-brokers-venues:1}
D'après l'\cref{dat:m1:exchanges-brokers-venues:groups}, calculer, pour le segment ICE
Exchanges, la part de chacun des postes suivants~: énergie~; actions au comptant et options sur actions~;
données et connectivité~; admissions à la cote.
\end{exercise}

\begin{exercise}[$\star$]\label{exo:m1:exchanges-brokers-venues:2}
Un contrat à terme a une valeur notionnelle de 280\,000~\$ et la \omterm{def:m1:exchanges-brokers-venues:exchange}{bourse} perçoit
0.70~\$ par côté. Exprimer les frais d'un aller-retour (achat puis vente) en points
de base du notionnel.
\end{exercise}

\begin{exercise}[$\star$]\label{exo:m1:exchanges-brokers-venues:3}
Pour chaque situation, nommer le modèle d'accès et préciser à qui incombent les contrôles
de risque pré-négociation~: (a) un gérant de fonds passe un ordre par téléphone à un \omterm{def:m1:the-sell-side:sales-trader}{sales-trader}~; (b)
l'algorithme d'un fonds quantitatif envoie ses ordres par la passerelle de son \omterm{def:m1:exchanges-brokers-venues:broker}{courtier}~;
(c) une société de trading colocalisée à la \omterm{def:m1:exchanges-brokers-venues:exchange}{bourse} utilise l'identifiant de son \omterm{def:m1:exchanges-brokers-venues:broker}{courtier}~;
(d) un \omterm{def:m1:what-a-trading-firm-does:market-maker}{teneur de marché} enregistré comme \omterm{def:m1:the-sell-side:sell-side}{broker-dealer}.
\end{exercise}

\begin{exercise}[$\star\star$]\label{exo:m1:exchanges-brokers-venues:4}
Une société négocie 60 millions d'actions par mois. Avec les frais de
l'\cref{ex:m1:exchanges-brokers-venues:crossover}, calculer le coût mensuel de
chaque modèle et l'économie réalisée par le meilleur par rapport au deuxième. À partir de quel
volume mensuel devient-il intéressant de demander à devenir membre~?
\end{exercise}

\begin{exercise}[$\star\star$]\label{exo:m1:exchanges-brokers-venues:5}
Un fournisseur d'\omterm{def:m1:exchanges-brokers-venues:dma}{accès sponsorisé} relève ses frais par action de 0.04 à
0.06~cent. Recalculer les deux seuils de bascule qui concernent l'\omterm{def:m1:exchanges-brokers-venues:dma}{accès sponsorisé}. Quel type
de client l'augmentation éloigne-t-elle~?
\end{exercise}

\begin{exercise}[$\star\star$]\label{exo:m1:exchanges-brokers-venues:6}
Avec les paramètres de l'\cref{ex:m1:exchanges-brokers-venues:five}, les
acteurs en place réduisent leurs frais et la nouvelle plateforme doit suivre~: $\nu$ tombe à
0.02~cent. Calculer la nouvelle part de marché au seuil de rentabilité et le profit à 10~\%.
\end{exercise}

\begin{exercise}[$\star\star\star$]\label{exo:m1:exchanges-brokers-venues:7}
\emph{Programmation.} Écrire \texttt{cheapest\_admissible(volume, max\_latency)}
en supposant pour chaque modèle d'accès une latence de 50 millisecondes, 2 millisecondes,
50 microsecondes et 40 microsecondes, respectivement. Une stratégie exige au plus
1 milliseconde et négocie 5 millions d'actions par mois~: quel modèle, à quel
coût mensuel, et combien la contrainte de latence lui coûte-t-elle par rapport au choix
sans contrainte~?
\end{exercise}

\begin{exercise}[$\star\star\star$]\label{exo:m1:exchanges-brokers-venues:8}
\emph{Trouver la faille.} Le plan d'une start-up affirme~: «~Les \omterm{def:m1:exchanges-brokers-venues:exchange}{bourses} en place facturent
0.30~cent par action pour prélever de la liquidité. Nous facturerons 0.05~cent. Les
traders minimisent leur coût, donc le volume passera chez nous.~» Repérer la grandeur que le plan a omise
dans le coût du trader, écrire la comparaison réelle que fait le trader, et
expliquer pourquoi la plateforme la moins chère peut être la plus coûteuse.
\end{exercise}

\section{Problème~: Construire une bourse}

\begin{problem}[{Problème du week-end --- plan d'affaires d'une nouvelle plateforme d'actions}]\label{pb:m1:exchanges-brokers-venues:1}
Un consortium de \omterm{def:m1:exchanges-brokers-venues:broker}{courtiers} prévoit une nouvelle plateforme pour un marché d'actions qui négocie 8
milliards d'actions par jour à un prix moyen de 45~\$, sur 252 jours.

\medskip\noindent\textbf{Partie I --- Recettes.}
\begin{enumerate}
\item La plateforme facture aux preneurs de liquidité 0.28~cent par action et verse aux fournisseurs de liquidité 0.25~cent.
Quels frais nets $\nu$ conserve-t-elle par action appariée~?
\item Exprimer $\nu$ en points de base de la valeur négociée.
\item Donner les recettes annuelles de transaction à 100~\% de part de marché, puis par
point de pourcentage de part.
\item Les recettes de données de marché seraient de 60~millions de dollars à 100~\% de part et
seraient proportionnelles à la part. Donner les recettes totales par point de part de marché.
\end{enumerate}

\medskip\noindent\textbf{Partie II --- Coûts et seuil de rentabilité.}
\begin{enumerate}[resume]
\item Les coûts technologiques s'élèvent à 18~millions de dollars par an, la réglementation et la surveillance à 9~
millions de dollars, le personnel à 11~millions de dollars. Donner la part de seuil de rentabilité $x^\star$.
\item Exprimer le seuil de rentabilité en actions par jour et en dollars négociés par jour.
\item Les membres du consortium exécutent ensemble 12~\% du volume du marché,
et pourraient acheminer un tiers de leurs ordres vers la nouvelle plateforme, où un appariement est
disponible. Supposer que les appariements sont disponibles pour la moitié de ces cas. Quelle part
cela donne-t-il~?
\item La plateforme est-elle rentable sur les seuls flux d'ordres des membres~? De combien~?
\end{enumerate}

\medskip\noindent\textbf{Partie III --- Acheter des parts de marché.}
\begin{enumerate}[resume]
\item Pour attirer les \omterm{def:m1:what-a-trading-firm-does:market-maker}{teneurs de marché}, la plateforme verse, pendant sa première année, une remise
de 0.32~cent tout en facturant encore 0.28~cent. Quelle est désormais $\nu$, et que coûte
une part de 4~\% en pertes de transaction~?
\item Ajouter les recettes de données et les coûts fixes~: quelle est la perte de la première année à 4~\%~?
\item La deuxième année, la plateforme revient à la grille tarifaire normale et détient 6~\%. Donner le
bénéfice, et le nombre de telles années nécessaires pour rembourser la perte de la première année.
\item Un acteur en place réagit en réduisant ses propres frais nets, et le nouvel entrant doit
s'aligner~: $\nu = 0.02$~cent. Recalculer $x^\star$.
\end{enumerate}

\medskip\noindent\textbf{Partie IV --- Jugement.}
\begin{enumerate}[resume]
\item Pourquoi les \omterm{def:m1:exchanges-brokers-venues:broker}{courtiers} veulent-ils cette plateforme même si elle se contente de l'équilibre~?
\item La plateforme envisage de devenir un \omterm{def:m1:exchanges-brokers-venues:ats}{système de négociation alternatif} (ATS) plutôt qu'une \omterm{def:m1:exchanges-brokers-venues:exchange}{bourse}. Citer deux
choses auxquelles elle renonce et une qu'elle économise.
\item Elle envisage d'introduire un délai d'une fraction de milliseconde sur les ordres entrants.
Quels traders cela attire-t-il, et lesquels cela repousse-t-il~?
\item Expliquer pourquoi la part de marché dans l'appariement obéit à l'économie d'un réseau,
et ce que cela implique pour le nombre de plateformes qui survivent.
\item Les recettes de données d'une \omterm{def:m1:exchanges-brokers-venues:exchange}{bourse} dépendent de sa part des \emph{cotations} aussi
bien que des transactions. Pourquoi cela pourrait-il récompenser l'affichage d'ordres
peu susceptibles d'être exécutés~?
\item Les clients en \omterm{def:m1:exchanges-brokers-venues:dma}{accès sponsorisé} se connectent directement à la plateforme. Que doit
la plateforme elle-même vérifier au sujet des \omterm{def:m1:exchanges-brokers-venues:broker}{courtiers} qui les parrainent~?
\item Énoncer le \emph{résultat nommé}~: la valeur négociée au seuil de rentabilité par jour, en
milliards de dollars, selon la grille tarifaire normale.
\item En une phrase~: s'agit-il d'une activité technologique ou d'une activité réglementaire~?
\end{enumerate}
\end{problem}

\section{Questions d'entretien}

\begin{interviewq}[$\star$ \iqroles{trader, développeur, chercheur}]\label{iq:m1:exchanges-brokers-venues:1}
Citez les moyens par lesquels une \omterm{def:m1:exchanges-brokers-venues:exchange}{bourse} gagne de l'argent. Lesquels comptent le plus pour
la structure de coûts d'une société de trading à haute fréquence~?
\end{interviewq}

\begin{interviewq}[$\star$ \iqroles{développeur, trader}]\label{iq:m1:exchanges-brokers-venues:2}
Quelle est la différence entre l'\omterm{def:m1:exchanges-brokers-venues:dma}{accès direct au marché} (DMA) et l'\omterm{def:m1:exchanges-brokers-venues:dma}{accès sponsorisé}~?
Qui est responsable si l'algorithme d'un client en \omterm{def:m1:exchanges-brokers-venues:dma}{accès sponsorisé} envoie dix mille
ordres erronés~?
\end{interviewq}

\begin{interviewq}[$\star\star$ \iqroles{développeur}]\label{iq:m1:exchanges-brokers-venues:3}
On vous demande d'implémenter des contrôles de risque pré-négociation que chaque ordre doit franchir.
Énumérez les contrôles et dites quelle contrainte leur budget de latence impose
à la conception.
\end{interviewq}

\begin{interviewq}[$\star\star$ \iqroles{chercheur, trader}]\label{iq:m1:exchanges-brokers-venues:4}
Pourquoi les \omterm{def:m1:exchanges-brokers-venues:exchange}{bourses} de contrats à terme exigent-elles des marges beaucoup plus élevées que les \omterm{def:m1:exchanges-brokers-venues:exchange}{bourses} d'actions~?
\end{interviewq}

\begin{interviewq}[$\star\star$ \iqroles{trader, chercheur}]\label{iq:m1:exchanges-brokers-venues:5}
Une nouvelle plateforme vous propose de vous verser 0.32~cent par action pour fournir de la liquidité alors que
la plateforme établie en verse 0.20. Votre taux d'exécution y est le cinquième de
celui de la plateforme établie. Comment décidez-vous où afficher vos cotations~?
\end{interviewq}

\begin{interviewq}[$\star\star\star$ \iqroles{chercheur, banque}]\label{iq:m1:exchanges-brokers-venues:6}
Les \omterm{def:m1:exchanges-brokers-venues:exchange}{bourses} vendent des \omterm{def:m1:exchanges-brokers-venues:feed}{flux de données de marché} plus rapides et la colocalisation à certains participants, et un
flux consolidé plus lent à tous les autres. Présentez l'argument le plus fort selon lequel
cette pratique est équitable, puis l'argument le plus fort selon lequel elle ne l'est pas.
\end{interviewq}
